\chapter{Independence of RVs, Functions of RVs \& PMF Convolution}

Independence is the bedrock of statistical modeling and machine learning assumptions---from Naive Bayes classifiers and i.i.d. training batches to random walk price modeling in quantitative finance. In this chapter, we master \textbf{Independence of Random Variables}, \index{Functions of Random Variables}\textbf{Functions of Random Variables}, \index{PMF Convolution}\textbf{PMF Convolution} for sums of independent RVs, and the exact probability laws governing \index{Extreme Values}\textbf{Extreme Values} ($\min(X, Y)$ and $\max(X, Y)$).

\vspace{0.5cm}
\begin{tcolorbox}[colback=primarylight, colframe=primarydark, title={\textbf{What You'll Learn in this Chapter}}]
\begin{itemize}
    \item Mathematical criterion for independence in bivariate and multivariate discrete distributions.
    \item Transformations and functions of random variables $Z = g(X, Y)$.
    \item The discrete convolution theorem for sums of independent random variables $Z = X + Y$.
    \item Order statistics and extreme value distributions: $V = \min(X, Y)$ and $W = \max(X, Y)$.
    \item Cognitive exam gotchas: Pairwise vs. Mutual independence and Uncorrelatedness vs. Independence.
\end{itemize}
\end{tcolorbox}

\section{Independence of Discrete Random Variables}

Two random variables $X$ and $Y$ are independent if learning the observed value of $Y$ provides zero information about the distribution of $X$.

\begin{definitionbox}{Independence Criterion}
\index{Independence Criterion}
Discrete random variables $X$ and $Y$ are \index{Independent}\textbf{Independent} (denoted $X \perp\!\!\!\perp Y$) if and only if their joint PMF factorizes into the product of their marginal PMFs for every pair of outcomes $(x, y) \in \mathcal{R}_X \times \mathcal{R}_Y$:
$$p_{X,Y}(x, y) = p_X(x) \cdot p_Y(y) \quad \forall (x, y)$$
Equivalently, for all $y$ where $p_Y(y) > 0$:
$$p_{X|Y}(x \mid y) = p_X(x) \quad \forall x$$
\end{definitionbox}

\begin{intuitionbox}{Checking Independence with Outer Products}
In linear algebra, an outer product $\mathbf{u} \mathbf{v}^T$ generates a rank-1 matrix where every row is a scalar multiple of every other row. The independence criterion $p_{X,Y}(x, y) = p_X(x) \cdot p_Y(y)$ establishes that a joint PMF table represents independent variables \textbf{if and only if the joint probability matrix has rank 1} and equals the outer product of its marginal vectors!
\end{intuitionbox}

\begin{lstlisting}[language=Python, caption=Data Science: Verifying Independence via Outer Product in NumPy]
import numpy as np

# Marginal vectors
p_X = np.array([0.4, 0.6]) # P(X=0), P(X=1)
p_Y = np.array([0.3, 0.7]) # P(Y=0), P(Y=1)

# Theoretical independent Joint PMF
joint_independent = np.outer(p_X, p_Y)
print("Theoretical Independent Joint PMF:\n", joint_independent)

# Empirical test: check if observed joint table matches outer product
observed_joint = np.array([[0.12, 0.28],
                           [0.18, 0.42]])
is_independent = np.allclose(observed_joint, joint_independent)
print(f"Are features independent? {is_independent}")
\end{lstlisting}

\begin{videocard}{IITM BS Lecture: Independence of Two Random Variables}{https://www.youtube.com/watch?v=YXaYJths2Ts}
Official lecture defining pairwise independence and testing factorization across joint PMF tables.
\end{videocard}

\begin{videocard}{IITM BS Lecture: Independence of Multiple Random Variables}{https://www.youtube.com/watch?v=Ndy0KTpmJAI}
Official lecture extending independence to $n$ random variables: $p_{X_1,\dots,X_n}(x_1,\dots,x_n) = \prod_{i=1}^n p_{X_i}(x_i)$.
\end{videocard}

\section{Functions of Random Variables and PMF Convolution}

If $X$ and $Y$ are random variables and $Z = g(X, Y)$ is a real-valued transformation, then $Z$ is itself a random variable. To compute the PMF of $Z$, we aggregate the joint probability mass over the pre-image of $z$:
$$p_Z(z) = \P(Z = z) = \sum_{\{(x, y) \colon g(x, y) = z\}} p_{X,Y}(x, y)$$

\begin{formulabox}{Discrete Convolution Formula for Sums}
When $X$ and $Y$ are \textbf{independent} discrete random variables, the PMF of their sum $Z = X + Y$ is given by the \index{Discrete Convolution}\textbf{Discrete Convolution} of $p_X$ and $p_Y$:
$$p_Z(z) = \P(X + Y = z) = \sum_{x \in \mathcal{R}_X} p_X(x) \cdot p_Y(z - x)$$
Equivalently, by symmetry of convolution: $p_Z(z) = \sum_{y \in \mathcal{R}_Y} p_X(z - y) \cdot p_Y(y)$.
\end{formulabox}

\begin{examplebox}{Convolution of Independent Poisson Variables}
Let $X \sim \text{Poisson}(\lambda_1)$ and $Y \sim \text{Poisson}(\lambda_2)$ be independent. The sum $Z = X + Y$ is also Poisson with parameter $\lambda_1 + \lambda_2$:
$$\P(X + Y = z) = \sum_{k=0}^z \left(\frac{e^{-\lambda_1}\lambda_1^k}{k!}\right) \left(\frac{e^{-\lambda_2}\lambda_2^{z-k}}{(z-k)!}\right) = \frac{e^{-(\lambda_1+\lambda_2)}}{z!} \sum_{k=0}^z \binom{z}{k} \lambda_1^k \lambda_2^{z-k} = \frac{e^{-(\lambda_1+\lambda_2)}(\lambda_1+\lambda_2)^z}{z!}$$
This additive closure is a fundamental property heavily exploited in queuing theory and event arrival modeling.
\end{examplebox}

\begin{videocard}{IITM BS Lecture: Functions of Random Variables \& PMF Convolution}{https://www.youtube.com/watch?v=I2S7nlCBkA8}
Official lecture computing PMFs of transformations $g(X,Y)$ and deriving discrete convolutions.
\end{videocard}

\section{Minimum and Maximum of Random Variables}

Extreme values $V = \min(X, Y)$ and $W = \max(X, Y)$ characterize bottleneck performance, fault tolerance, and multi-sensor reliability.

\begin{formulabox}{Min and Max CDF Formulas for Independent RVs}
For mutually independent random variables $X$ and $Y$:
\begin{itemize}
    \item \textbf{Maximum $W = \max(X, Y)$ (CDF factorizes directly):}
    $$F_W(w) = \P(\max(X, Y) \le w) = \P(X \le w \text{ and } Y \le w) = F_X(w) \cdot F_Y(w)$$
    \item \textbf{Minimum $V = \min(X, Y)$ (Survival function factorizes):}
    $$S_V(v) = \P(\min(X, Y) > v) = \P(X > v \text{ and } Y > v) = (1 - F_X(v)) \cdot (1 - F_Y(v))$$
    Therefore, the CDF of the minimum is:
    $$F_V(v) = 1 - S_V(v) = 1 - [1 - F_X(v)][1 - F_Y(v)]$$
\end{itemize}
\end{formulabox}

\begin{videocard}{IITM BS Lecture: Minimum \& Maximum of Two Random Variables}{https://www.youtube.com/watch?v=sgzr8rTaAcI}
Official lecture deriving probability mass functions and cumulative distributions for $\min(X,Y)$ and $\max(X,Y)$.
\end{videocard}

\section{Exam Gotchas \& Cognitive Traps}

\begin{trapbox}{Pairwise Independence Does NOT Imply Mutual Independence}
\textbf{The Trap:} Assuming that if every pair of variables satisfies $p_{X_i, X_j} = p_{X_i} p_{X_j}$, then the entire collection is mutually independent.

\textbf{The Classic Counterexample (Bernstein's Triplet):}
Let $X \sim \text{Bernoulli}(0.5)$ and $Y \sim \text{Bernoulli}(0.5)$ be fair independent coin flips. Define $Z = X \oplus Y = (X + Y) \bmod 2$.
\begin{itemize}
    \item $X \perp\!\!\!\perp Y$, $Y \perp\!\!\!\perp Z$, and $X \perp\!\!\!\perp Z$ are all \textbf{pairwise independent}!
    \item Yet they are \textbf{mutually dependent}: knowing $X$ and $Y$ completely determines $Z$ with certainty! In fact, $\P(X=1, Y=1, Z=1) = 0 \neq p_X(1)p_Y(1)p_Z(1) = \frac{1}{8}$.
\end{itemize}
\end{trapbox}

\begin{trapbox}{Convolution Requires Independence}
\textbf{The Trap:} Applying the convolution formula $p_{X+Y}(z) = \sum_x p_X(x) p_Y(z-x)$ when $X$ and $Y$ are correlated or dependent.

\textbf{The Reality:} If $X$ and $Y$ are dependent, you \textbf{must} sum the true joint PMF:
$$p_{X+Y}(z) = \sum_x p_{X,Y}(x, z - x)$$
Multiplying marginals $p_X(x) p_Y(z-x)$ is strictly invalid whenever dependence exists.
\end{trapbox}

\section{Comparative Synthesis Matrix}

\begin{table}[htbp]
\centering
\small
\begin{tabular}{p{3.2cm}p{3.8cm}p{3.8cm}p{3.5cm}}
\toprule
\textbf{Property} & \textbf{Statistical Independence} & \textbf{Uncorrelatedness} & \textbf{Mutual Exclusivity (Disjoint)} \\
\midrule
\textbf{Definition} & $p_{X,Y}(x,y) = p_X(x)p_Y(y)$ & $\Cov(X, Y) = 0$ & $\P(X=x, Y=y) = 0$ (for $x \neq y$) \\
\addlinespace
\textbf{Strength} & Strongest condition (all moments) & Linear dependence only & Extreme negative dependence! \\
\addlinespace
\textbf{Implication Direction} & Independence $\implies$ Uncorrelated & Uncorrelated $\not\implies$ Independence & Disjoint with non-zero prob $\implies$ \textbf{Never} independent! \\
\addlinespace
\textbf{Joint Support Shape} & Must be Cartesian product grid $\mathcal{R}_X \times \mathcal{R}_Y$ & Can take arbitrary shapes & Strictly off-diagonal points \\
\addlinespace
\textbf{Data Science Impact} & Enables feature product rules (Naive Bayes) & Features have zero Pearson linear correlation & One-hot encoded categorical indicators \\
\bottomrule
\end{tabular}
\caption{Contrasting Statistical Independence, Uncorrelatedness, and Mutual Exclusivity.}
\end{table}

\newpage
\section{Practice Exercises}
\textit{Detailed step-by-step solutions are available in the Appendix.}

\subsection*{Part A: Analytical \& Computational Problems}

\begin{enumerate}
\item \textbf{Testing Independence from Joint PMF:} Let $X \in \{0, 1\}$ and $Y \in \{0, 1\}$ have joint PMF:
    $$p_{X,Y}(0,0) = 0.4, \quad p_{X,Y}(0,1) = 0.2, \quad p_{X,Y}(1,0) = 0.2, \quad p_{X,Y}(1,1) = 0.2$$
    Are $X$ and $Y$ independent? Justify by comparing $p_{X,Y}(0,0)$ with $p_X(0) \cdot p_Y(0)$.

\item \textbf{PMF Convolution of Independent Uniforms:} Let $X$ and $Y$ be independent discrete uniform random variables on $\{1, 2, 3\}$. Find the PMF of $Z = X + Y$.

\item \textbf{Sum of Independent Poisson RVs:} Let $X \sim \text{Poisson}(\lambda_1 = 2)$ and $Y \sim \text{Poisson}(\lambda_2 = 3)$ be independent. State the distribution of $Z = X + Y$ and compute $\P(Z = 1)$.

\item \textbf{Maximum of Two Geometric RVs:} Let $X$ and $Y$ be independent geometric random variables with success probability $p = 0.5$ ($\P(X \le k) = 1 - (1-p)^k$). Find the CDF $F_W(2)$ for $W = \max(X, Y)$.

\item \textbf{Data Science Application (Server Latency Redundancy Minimum):} A high-availability web app routes queries to two parallel microservices. Let $X, Y \in \{1, 2, 3\}$ seconds be independent service response times with $p_X(k) = p_Y(k) = 1/3$. Compute the PMF of the overall system response time $V = \min(X, Y)$.
\end{enumerate}

\subsection*{Part B: Conceptual Diagnostic Traps (Exam Gotchas)}

\begin{enumerate}[resume]
\item \textbf{Diagnostic Trap (Pairwise vs. Mutual Independence):}
Let $A, B \stackrel{iid}{\sim} \text{Bernoulli}(0.5)$ and define $C = (A + B) \bmod 2$. Which of the following statements is strictly correct?
\begin{enumerate}
    \item $A, B,$ and $C$ are mutually independent because $\P(A=1, B=1) = \P(A=1)\P(B=1)$.
    \item Any pair chosen from $\{A, B, C\}$ is independent, but the triplet $\{A, B, C\}$ is mutually dependent.
    \item $C$ is independent of $A$ only if $B = 0$.
    \item Since $C$ is a deterministic function of $A$ and $B$, $A$ and $C$ cannot be pairwise independent.
\end{enumerate}

\item \textbf{Diagnostic Trap (Support Geometry and Independence):}
A bivariate distribution $p_{X,Y}(x, y)$ has non-zero probability mass strictly inside the triangular domain $\mathcal{T} = \{(x, y) \in \mathbb{N}^2 \colon 1 \le x \le y \le 4\}$. Can $X$ and $Y$ be independent?
\begin{enumerate}
    \item Yes, provided that $p_{X,Y}(x, y)$ is uniform across all points in $\mathcal{T}$.
    \item Yes, if the marginal distributions $p_X$ and $p_Y$ have identical variances.
    \item No, because an independent joint distribution must have a rectangular Cartesian product support $\mathcal{R}_X \times \mathcal{R}_Y$.
    \item No, because discrete random variables can never be independent on finite supports.
\end{enumerate}
\end{enumerate}
